\documentclass[11pt]{article}
\usepackage{mathblatt}
% gesetzt von werkzeuge/setzer.py; Lösungen nach bau/bauregeln.md „Lösungen“.
\newcommand{\kennung}{FA5}
\begin{document}
\pfheftstil
\fancyfoot[C]{{\footnotesize\color{mbgrau}\kennung\hfill\thepage}}
\pfheftkopf{Zuordnungstypen erkennen und anwenden}{Klasse 7 \textperiodcentered{} Lösungen}

\begin{pfloesung}{}
\lz{1b)}{proportional}{$6 : 2 = 12 : 4 = 15 : 5 = 24 : 8 = 3$}{}
\end{pfloesung}
\begin{pfloesung}{}
\lz{2.}{antiproportional}{Quotienten $30$, $7{,}5$, … nicht gleich; Produkte $1 \cdot 30 = 2 \cdot 15 = 3 \cdot 10 = 5 \cdot 6 = 30$}{}
\end{pfloesung}
\begin{pfloesung}{}
\lz{3.}{keins von beiden}{Quotienten $5 : 1 = 5$, $7 : 2 = 3{,}5$; Produkte $1 \cdot 5 = 5$, $2 \cdot 7 = 14$; beide nicht gleich}{}
\end{pfloesung}
\begin{pfloesung}{}
\lz{4.}{(a) keins \ (b) proportional \ (c) antiproportional}{(a) Gerade, aber nicht durch den Ursprung; (b) Gerade durch den Ursprung; (c) fallende Kurve}{}
\end{pfloesung}
\begin{pfloesung}{}
\lz{5.}{keins von beiden; kein Dreisatz}{Keins von beiden; Quotienten $9 : 2 = 4{,}5$, $13 : 4 = 3{,}25$ – nicht gleich; Produkte $2 \cdot 9 = 18$, $4 \cdot 13 = 52$ – nicht gleich. Also kein Dreisatz (der Verleih nimmt eine Grundgebühr von 5\,€ und 2\,€ je Stunde).}{}
\end{pfloesung}
\begin{pfloesung}{}
\lz{6b)}{3{,}20\,€}{proportional: $2 : 5 = 0{,}40$; $0{,}40 \cdot 8 = 3{,}20$\,€}{}
\end{pfloesung}
\begin{pfloesung}{}
\lz{7.}{6 Stunden}{antiproportional: $3 \cdot 8 = 24$; $24 : 4 = 6$\,h}{}
\end{pfloesung}
\begin{pfloesung}{}
\lz{8.}{16 Minuten}{antiproportional (doppelt so schnell, halb so lange): $12 \cdot 20 = 240$; $240 : 15 = 16$\,min}{}
\end{pfloesung}
\begin{pfloesung}{}
\lz{9.}{33 Minuten}{proportional (doppelt so weit, doppelt so lange): $22 : 4 = 5{,}5$; $5{,}5 \cdot 6 = 33$\,min}{}
\end{pfloesung}
\begin{pfloesung}{}
\lz{10.}{Nein, ihr Anteil ist 105\,€; es fehlen 5\,€.}{Nein; antiproportional: $6 \cdot 140 = 840$; $840 : 8 = 105$\,€; $105 - 100 = 5$\,€ fehlen.}{}
\end{pfloesung}
\begin{pfloesung}{}
\lz{11b)}{4{,}20\,€}{$12 \cdot 35 = 420$\,ct $= 4{,}20$\,€}{}
\end{pfloesung}
\begin{pfloesung}{}
\lz{12.}{6 Minuten}{$900 : 2{,}5 = 360$\,s; $360 : 60 = 6$\,min}{}
\end{pfloesung}
\begin{pfloesung}{}
\lz{13.}{16\,km/h}{$15$\,min $\cdot 4 = 60$\,min; $4$\,km $\cdot 4 = 16$\,km in 1 Stunde}{}
\end{pfloesung}
\begin{pfloesung}{}
\lz{14.}{45 Minuten}{$18 : 24 = 0{,}75$\,h; $0{,}75 \cdot 60 = 45$\,min}{}
\end{pfloesung}
\begin{pfloesung}{}
\lz{15.}{Das stimmt: 324\,€ gespart.}{$12\,000 : 100 = 120$; $120 \cdot 6 = 720$\,l; $120 \cdot 4{,}5 = 540$\,l; $720 - 540 = 180$\,l; $180 \cdot 1{,}80 = 324$\,€ $> 300$\,€}{}
\end{pfloesung}
\begin{pfloesung}{}
\lz{16.}{keins von beiden}{Quotienten $0{,}7$, $0{,}45$; Produkte $70$, $180$; beide nicht gleich}{}
\end{pfloesung}
\begin{pfloesung}{}
\lz{17.}{8 Minuten}{proportional: $45 : 3 = 15$ Seiten je Minute; $120 : 15 = 8$\,min}{}
\end{pfloesung}
\begin{pfloesung}{}
\lz{18.}{10 Tage}{antiproportional: $8 \cdot 15 = 120$; $120 : 12 = 10$ Tage}{}
\end{pfloesung}
\begin{pfloesung}{}
\lz{19.}{2\,h 15\,min}{$5400 : 40 = 135$\,min; $135 - 120 = 15$, also 2\,h 15\,min}{}
\end{pfloesung}
\begin{pfloesung}{}
\lz{20.}{25 Personen: Bahn (Bus 24\,€); 30 Personen: Bus (20\,€)}{Bei 25: Bahn; Bus $600 : 25 = 24$\,€ $> 22$\,€. Bei 30: Bus; $600 : 30 = 20$\,€ $< 22$\,€.}{}
\end{pfloesung}
\end{document}
